A introdução mostrou que os agentes de IA já são usados para modernizar sistemas legados, mas que se sabe pouco sobre como é usá-los no dia a dia de quem desenvolve, em um sistema real e de tamanho pequeno. Este capítulo diz o que o trabalho quer alcançar para tratar dessa lacuna e mostra em qual capítulo cada parte é feita.

O objetivo geral é identificar os desafios para converter e atualizar sistemas legados para novas tecnologias por meio de agentes de IA, do ponto de vista do desenvolvimento. Aqui, o ponto de vista do desenvolvimento é o de quem escreve, mantém e confere o código, e não o de quem administra ou usa o sistema. Um desafio, neste trabalho, é uma dificuldade que aparece quando o agente faz esse serviço e que o desenvolvedor precisa perceber e tratar, como uma regra do sistema que o agente não percebeu ou um resultado difícil de conferir. Esse objetivo responde à pergunta apresentada na introdução.

Para alcançá-lo, o trabalho olha para o assunto de dois lados: o que a literatura já relata e o que acontece na prática. Por isso, o objetivo geral é dividido em quatro etapas, que são os objetivos específicos, e cada uma é feita em um capítulo.

A primeira etapa é reunir o que estudos recentes relatam e organizar os desafios encontrados. O Capítulo~3, de fundamentação teórica, explica como esses estudos foram buscados e escolhidos e depois apresenta o assunto em quatro partes: como a modernização aparece nos estudos, como os agentes de IA são usados nessa tarefa, como o resultado é verificado e quais desafios a literatura relata. Esses desafios são reunidos em um quadro, que serve de lista do que deve ser observado no estudo de caso e que também mostra o que ainda falta ser estudado.

A segunda etapa é definir como o estudo de caso é feito. O Capítulo~4 explica o tipo de pesquisa e apresenta o sistema estudado, feito no Microsoft Power Apps: para que ele serve, como é organizado e quais regras dele foram escolhidas, de complexidades diferentes. Depois mostra que os dados usados são inventados, como os resultados do sistema original foram obtidos em um app de teste, como o agente de IA foi usado e como o resultado do código novo é comparado com o esperado. Assim, o leitor entende o que foi feito e consegue repetir o método.

A terceira etapa é converter essas regras para Python, com a ajuda de um agente de IA, e conferir o resultado por testes automáticos, que são pequenos programas que comparam o que o código devolve com o que se espera. O Capítulo~5 mostra o que foi convertido e o resultado dos testes.

A quarta etapa é observar os problemas que apareceram nessa conversão, ligá-los, quando possível, aos desafios do quadro da literatura e sugerir cuidados ao desenvolvedor. Esta parte também está no Capítulo~5, junto com os limites do estudo. O Capítulo~6 fecha o trabalho: responde à pergunta da introdução, retoma os limites e aponta o que pode ser feito em trabalhos futuros.

Para que o trabalho possa ser acompanhado de perto, ele faz algumas escolhas de escopo. O foco é o que o desenvolvedor encontra ao pedir a um agente de IA que atualize um sistema, e questões de gestão, de uso pelos usuários e de custo ficam de fora. Em vez de converter um sistema inteiro, o estudo de caso trata de quatro regras de um sistema feito no Microsoft Power Apps, de complexidades diferentes, que foram convertidas para Python. Também foi usado um só agente de IA, o que permite acompanhar o processo de perto, mas não permite comparar agentes entre si. A conferência do resultado é feita por testes automáticos, e a qualidade do código além disso, como desempenho ou organização, não é avaliada. Os limites que essas escolhas trazem para os resultados são discutidos no Capítulo~5.

A contribuição do trabalho é de três tipos. Primeiro, ele organiza em um quadro os desafios que a literatura relata, o que ajuda a enxergar, de uma vez só, o que pode dar errado nesse tipo de trabalho. Segundo, mostra em um caso pequeno e concreto como alguns desses desafios aparecem na prática e quais não puderam ser observados. Terceiro, deixa um método simples, com dados inventados, um app de teste e testes automáticos, que outro desenvolvedor pode repetir em um sistema seu.
